\documentclass[UTF8,openany,11pt]{ctexbook}
\usepackage[a4paper,margin=2.2cm]{geometry}
\usepackage{booktabs,longtable,array}
\usepackage{hyperref}
\usepackage{xcolor}
\usepackage{enumitem}
\usepackage{fancyvrb}
\usepackage{titlesec}
\usepackage{setspace}
\usepackage{tcolorbox}
\hypersetup{colorlinks=true,linkcolor=blue,urlcolor=blue}
\setlist{nosep,leftmargin=2em}
\setstretch{1.18}
\sloppy
\titleformat{\chapter}[display]{\bfseries\Huge}{第\thechapter 章}{0.4em}{}
\newtcolorbox{keypoint}{colback=blue!4!white,colframe=blue!45!black,title=核心结论}
\newtcolorbox{warningbox}{colback=orange!6!white,colframe=orange!60!black,title=常见误解}
\newcommand{\pathref}[1]{\texttt{#1}}
\newcommand{\term}[1]{\textbf{#1}}
\begin{document}

\frontmatter
\begin{titlepage}
\centering
\vspace*{2.5cm}
{\Huge\bfseries 系统理解 Pi\par}
\vspace{0.4cm}
{\Large 从终端 Coding Agent 到 Agent Harness\par}
\vspace{1.4cm}
{\large 核心主线教材 v0.1\par}
\vfill
\begin{tcolorbox}[colback=gray!5!white,colframe=gray!50!black,width=0.88\textwidth]
\textbf{源码基线}\quad earendil-works/pi\\
\textbf{Commit}\quad \texttt{dcfe36c79702ec240b146c45f167ab75ecddd205}\\
\textbf{读者目标}\quad 能运行、追踪、解释并开始修改 Pi；为后续源码深读和 mini Pi 复刻建立框架。
\end{tcolorbox}
\vspace{1cm}
{\large 2026\par}
\end{titlepage}

\chapter*{如何使用这本书}
\addcontentsline{toc}{chapter}{如何使用这本书}
这不是一份命令列表，也不是把 Pi 的文件逐行翻译成中文。Pi 是一个正在快速演进的 TypeScript monorepo；若从最大的 UI 文件开始，或者从四千多个 commit 线性阅读，几乎一定会失去方向。本书以一个问题组织：\textbf{用户在终端输入一条指令之后，系统怎样把它变成可靠、可恢复、可扩展的 Agent 行为？}

每章遵循相同节奏：
\begin{enumerate}
\item 先建立一个可以观察的场景；
\item 再锁定少量源码锚点；
\item 解释数据、状态和控制如何流动；
\item 给出检查问题，确认你不是只“看过”代码。
\end{enumerate}

请不要一口气读完。建议每次只读一章，运行一个命令或一个 faux-provider 测试，然后在自己的笔记中写下：场景、入口、状态所有者、协议、失败路径、仍不理解的问题。

\tableofcontents

\mainmatter

\chapter{Pi 是什么，以及它不是什么}

\section{一句话定位}
Pi 是一个 \term{Agent harness}：它为语言模型提供消息协议、流式响应、工具调用、会话、上下文管理、多 Provider 适配和终端交互外壳。其最知名的产品是终端 Coding Agent，但仓库并不等于“一个调用模型的 CLI”。

根目录 \pathref{README.md} 列出四个核心包：
\begin{longtable}{p{0.27\textwidth}p{0.26\textwidth}p{0.37\textwidth}}
\toprule
包 & 职责 & 不应该负责什么\\
\midrule
\texttt{packages/ai} & Provider、模型目录、认证、协议适配、流 & 文件编辑、终端 UI\\
\texttt{packages/agent} & Agent state、turn、tool loop、事件 & 具体 coding tool、主题\\
\texttt{packages/coding-agent} & CLI、工具、session、扩展、产品策略 & 重写 Agent loop\\
\texttt{packages/tui} & 输入编辑、组件、ANSI 差量渲染 & 模型调用、会话策略\\
\bottomrule
\end{longtable}

\begin{keypoint}
Pi 的关键不是某一个 prompt，而是边界。模型层不知道文件系统；Agent core 不知道主题和键盘；UI 消费事件而不拥有第二份对话状态。只要边界保持，任何一层都可以替换。
\end{keypoint}

\section{为什么 Agent 必须是循环}
纯聊天近似一次请求：
\begin{Verbatim}[fontsize=\small]
user message -> LLM -> text
\end{Verbatim}
Coding Agent 则是闭环：
\begin{Verbatim}[fontsize=\small]
user -> LLM -> tool call -> host executes read/bash/edit
                 ^                    |
                 |---- tool result ---+
\end{Verbatim}
模型不会直接读磁盘。它只能生成 ``调用 read，参数为 path'' 的结构化意图；Pi 负责验证、执行、记录结果，再把结果作为消息送回模型。一次用户请求可以包含多个 turn：每个 turn 有一次模型响应，可能跟着一批工具执行。

\section{Pi 与 LangChain、LangGraph}
Pi 当前仓库没有 LangChain 依赖。它和 LangChain/LangGraph 都处理模型、tools、memory、callback，但抽象中心不同：Pi 把\term{单 Agent 的消息--工具循环}作为核心原语；LangGraph 更适合将业务显式建模为多个节点和边的状态图。Pi 适合学习交互式 coding agent 的运行时；若将来要做多 Agent 审批流、复杂工作流编排，再补 LangGraph 更合理。

\section{本书的完成定义}
读完核心版，你应该能：
\begin{itemize}
\item 解释四层包如何协作；
\item 追踪一次 prompt 从 CLI 到 tool result 再到最终回答；
\item 运行无需 API key 的 faux-provider 测试；
\item 知道 session、context、Provider cache 与模型内部 KV cache 的边界；
\item 选择一个低风险行为做出最小改动，并找到对应测试。
\end{itemize}

\chapter{环境、运行方式与第一条验证锚点}

\section{建立可重复环境}
Pi 当前要求 Node 22.19 或更高版本。请在完整 clone 根目录执行：
\begin{Verbatim}[fontsize=\small]
cd D:\projects\pi-source-analysis
npm install --ignore-scripts
.\pi-test.ps1 --help
.\pi-test.ps1 --list-models
\end{Verbatim}

\pathref{pi-test.ps1} 是 Windows 的源码启动包装：它定位根目录 \texttt{tsx.cmd}，然后执行 \pathref{packages/coding-agent/src/cli.ts}。它可附加 \texttt{--no-env}，清除常见 Provider 凭据，避免实验意外使用真实账号。

\section{三种运行模式}
在 \pathref{packages/coding-agent/src/main.ts} 中，\texttt{resolveAppMode(parsed, stdinIsTTY, stdoutIsTTY)} 的规则非常重要：
\begin{Verbatim}[fontsize=\small]
--mode rpc                => rpc
--mode json               => json
--print 或任一流非 TTY  => print
其他                      => interactive
\end{Verbatim}
这证明终端 UI 不是 Agent 的前提。相同的 session/runtime 可以被人类终端、管道脚本和 RPC 客户端复用。

\section{为什么 \texttt{--list-models} 可能没有模型}
\texttt{--list-models} 并非简单列出一个硬编码数组。它经过模型注册表和认证可用性判断。没有 API key 或 OAuth 账号时，Pi 正确地提示 ``No models available'' 并指向 provider/model 文档。它不是安装失败，而是模型层刻意不把“已知模型”误报为“当前可请求模型”。

\section{从 faux provider 开始}
不要把第一次端到端学习建立在真实 API 上。Pi 自带 faux provider，可以确定性地输出指定的 assistant message、thinking、tool call、延迟与 usage。

\begin{Verbatim}[fontsize=\small]
cd packages\agent
.\node_modules\.bin\vitest.cmd --run test\e2e.test.ts `
  -t "executes tools and tracks pending tool calls"
\end{Verbatim}

该测试位于 \pathref{packages/agent/test/e2e.test.ts}。它让 faux 模型先返回一个 \texttt{calculate} tool call，再返回包含计算结果的文本；测试断言：
\begin{itemize}
\item transcript 中存在 tool result；
\item 最后消息是 assistant 且含正确答案；
\item \texttt{pendingToolCalls} 在 \texttt{tool\_execution\_start} 时包含 call ID，在 end 时为空。
\end{itemize}

\begin{warningbox}
不要把 ``测试通过'' 理解为 ``理解完成''。它只是你的第一个固定观测点：以后读 loop、加工具、研究并发时都可以回到这条测试验证基本闭环没有被破坏。
\end{warningbox}

\chapter{从终端输入到 AgentSession}

\section{薄入口的价值}
\pathref{packages/coding-agent/src/cli.ts} 很短：设定进程名，标记 \texttt{PI\_CODING\_AGENT}，配置 HTTP dispatcher，然后调用 \texttt{main(process.argv.slice(2))}。这种薄入口避免 import 副作用、参数解析、网络配置、业务状态和 UI 互相纠缠。

\section{\texttt{main.ts} 是产品装配器}
\pathref{main.ts} 不是 Agent loop。它依次处理参数、stdin/附件、session 选择、settings、project trust、resource loader、认证和模型，然后创建 runtime 并分派 mode。你读它时只回答四个问题：
\begin{enumerate}
\item 什么条件决定 mode？
\item 哪个对象负责当前 cwd？
\item session 从何处恢复或创建？
\item Agent、工具、模型和 extensions 在哪里装配？
\end{enumerate}

\section{为什么需要 AgentSession}
\pathref{packages/agent/src/agent.ts} 的 \texttt{Agent} 是通用 runtime：它维护消息、工具、streaming 状态和队列。Coding Agent 还需产品能力：JSONL session、模型选择、compaction、bash、project resources、extension lifecycle。于是 \pathref{packages/coding-agent/src/core/agent-session.ts} 包装 core Agent，提供 coding-agent 专有策略。

\begin{Verbatim}[fontsize=\small]
CLI / print / RPC / interactive
              |
       AgentSessionRuntime
              |
         AgentSession
              |
            Agent
              |
      pi-ai Model / Provider
\end{Verbatim}

\section{cwd 是服务边界}
\pathref{agent-session-runtime.ts} 在 session 切换时不会只修改一个 cwd 字符串。它会关闭旧 session 生命周期、让 host 解除旧 UI/extension 绑定、按新 cwd 重建 services，然后绑定新 session。原因是 settings、skills、extensions、context files 和 trust 都可能依赖项目目录。复用旧项目资源会构成权限和一致性问题。

\chapter{消息、流与状态：Agent 的数据模型}

\section{两种消息层}
Pi 区分 \term{AgentMessage} 与 LLM 可理解的 \term{Message}。前者可以包含应用专用记录；后者只能是 Provider 支持的 user、assistant、toolResult 等标准消息。边界是：
\begin{Verbatim}[fontsize=\small]
AgentMessage[]
  -> transformContext()  // 裁剪、注入、压缩
  -> convertToLlm()     // 过滤或转换自定义消息
  -> Message[]
  -> Provider
\end{Verbatim}

这条边界避免 UI 通知、分支摘要或内部 session entry 意外泄漏给模型，同时允许产品在调用模型前重写上下文。

\section{partial 与 final}
Provider stream 可能产生 start、text delta、thinking delta、tool-call delta、done/error。用户希望马上看到 token；session 却只能保存完整、稳定的消息。因此 \pathref{packages/agent/src/agent-loop.ts} 的 \texttt{streamAssistantResponse()} 在 start 时持有 partial assistant，在每个 delta 更新它，最终通过 \texttt{response.result()} 获取 final assistant message。

\begin{keypoint}
partial message 是渲染状态，不是持久事实。若把 partial 直接写入 JSONL 或作为下一轮完整上下文，工具参数和 usage 都可能不完整。
\end{keypoint}

\section{Agent event 不是日志}
典型事件序列：
\begin{Verbatim}[fontsize=\small]
agent_start
turn_start
message_start(user) / message_end(user)
message_start(assistant)
message_update(...)
message_end(assistant)
tool_execution_start / update / end
message_start(toolResult) / message_end(toolResult)
turn_end
agent_end
\end{Verbatim}
\texttt{Agent.processEvents()} 的顺序是：先归约内部 state，后按注册顺序 await listener。这个顺序让 session persistence、UI 和 tool hook 观察同一份已更新的状态，而不是竞争谁先看到 assistant message。

\chapter{Agent loop：模型、工具与下一轮}

\section{先看最小伪代码}
\begin{Verbatim}[fontsize=\small]
append user prompt
while true:
  assistant = stream model(context)
  append final assistant
  calls = assistant.toolCalls
  results = execute calls
  append results
  if no calls: stop
\end{Verbatim}
真实 Pi 在此基础上增加 context transform、abort、错误处理、parallel/sequential batch、steering、follow-up、hooks、dynamic next-turn state 和 compaction 协作。

\section{源码阅读顺序}
打开 \pathref{packages/agent/src/agent-loop.ts}，不要从底部 helper 开始。依次读：
\begin{enumerate}
\item \texttt{runAgentLoop()}：新 prompt 如何进入 context；
\item \texttt{runLoop()}：双层循环如何组织；
\item \texttt{streamAssistantResponse()}：唯一的模型调用边界；
\item \texttt{executeToolCalls()}：如何选择串行或并行；
\item tool prepare/execute/finalize：验证、hook、执行、错误和结果。
\end{enumerate}

\section{为什么有两层循环}
内层处理 ``还有 tool call'' 或 ``有 steering message''；外层在本应结束时查询 follow-up queue。语义不同：
\begin{itemize}
\item \term{abort}：取消正在进行的流或工具；
\item \term{steering}：当前 tool batch 完成后优先插入的新指令；
\item \term{follow-up}：Agent 原本完成后才开始的新工作。
\end{itemize}
把三者都实现为 ``追加 prompt'' 会导致 tool 未完成、历史顺序混乱或用户意图被错误解释。

\section{工具批次的顺序规则}
Pi 默认可并行执行工具，但 preflight 按顺序完成；若任意工具要求 \texttt{executionMode: "sequential"}，整批改为串行。并发时 UI 的 \texttt{tool\_execution\_end} 可以按完成顺序即时出现；而持久化 tool result 仍按 assistant 原始 source order 写入 transcript。

这两个顺序同时存在是成熟实现常见的细节：人想尽快看到完成；模型和 session 需要稳定、可重放的历史。

\section{两个必须理解的失败分支}
\begin{enumerate}
\item assistant 因输出长度停止且带 tool call：参数可能截断，即使 JSON 恰好可解析也不安全；Pi 拒绝执行并要求模型重发。
\item 工具抛异常：Agent 转成 \texttt{isError: true} 的 tool result，让模型有机会修复调用；不能静默吞错，也不应让整个 session 崩溃。
\end{enumerate}

\chapter{Coding tools：文件、命令与安全边界}

\section{工具不只是函数}
一个 Coding Agent tool 至少有模型可读描述、参数 schema、执行器、可展示结果、details、错误语义、取消支持和输出边界。Coding Agent 的内建工具在 \pathref{packages/coding-agent/src/core/tools/}；装配入口在 \pathref{core/tools/index.ts}。

\section{read、edit、write、bash}
\begin{longtable}{p{0.18\textwidth}p{0.35\textwidth}p{0.37\textwidth}}
\toprule
工具 & 主要目的 & 主要风险\\
\midrule
read & 获取文件上下文 & 大文件、二进制、路径出界、token 爆炸\\
edit & 受控 find/replace & 旧内容不匹配、重复替换、并发修改\\
write & 创建/覆盖内容 & 误覆盖、目录边界、部分写入\\
bash & 执行命令 & shell 注入、挂起、子进程、凭据和网络\\
\bottomrule
\end{longtable}

\section{截断是上下文工程}
模型上下文有限。一次巨大命令输出或读取依赖目录，可能挤掉用户任务和系统提示词。Pi 在工具返回层截断，而非等 Provider 报 context overflow；并把截断事实明确放入返回内容。这样模型、UI、session 对看到的结果一致。

\section{Pi 默认不是沙箱}
Pi README 明确说明：它以启动进程的用户权限访问文件、进程、网络和凭据，没有内建权限系统。Prompt 不是权限系统；让模型“不要删除文件”不能限制 shell。需要更强边界时，使用容器、VM、OpenShell 或 Gondolin 等外部方案。学习工具实现时，请始终把``路径检查''、``输出限制''与``真正隔离''视为三个不同层次。

\chapter{会话、上下文与压缩}

\section{transcript 与模型 context 不等同}
Pi 可以在 session 文件中保存 user/assistant/tool result 之外的分支、摘要、custom entry；模型每轮看到的却是经过 transform 和 conversion 的上下文。这样做的目的不是隐藏历史，而是让可恢复记录、UI 呈现和有限 context window 可以分别优化。

\section{为什么选择 JSONL}
\pathref{packages/coding-agent/src/core/session-manager.ts} 采用逐行 JSON：
\begin{itemize}
\item 新 entry 可追加，不必重写完整对话；
\item 进程中断后，已完成行仍可读取；
\item 导出、分支和诊断能以 entry 为单位处理；
\item header、summary、custom entry 可自然共存。
\end{itemize}
恢复时不应静默跳过坏行；若历史被截断，系统必须把错误位置和恢复策略明确给用户。

\section{compaction 不是删除}
当上下文接近模型窗口，Pi 的 compaction 选择一段历史生成 summary，并以显式 entry 取代模型上下文中的旧细节。正确的 compaction 要满足：
\begin{enumerate}
\item 用户能知道摘要发生过；
\item summary 有清晰切点和来源；
\item 失败或 abort 不会损坏原 session；
\item 后续消息、queue 和 retry 的顺序保持一致。
\end{enumerate}
\pathref{packages/coding-agent/src/core/compaction/} 是后续深读专题；本阶段先牢牢记住：``保存了什么''、``显示了什么''、``发给模型什么'' 可能不同，且差异必须可解释。

\chapter{Provider、prompt cache 与 KV cache}

\section{Models 与 Provider}
\pathref{packages/ai/src/models.ts} 把职责拆为：
\begin{itemize}
\item \texttt{Provider}：自己的 model list、认证语义、请求流；
\item \texttt{Models}：查找 provider/model、解析认证、合并请求选项、委托 stream；
\item adapter：把 OpenAI、Anthropic、Google 等不同协议归一为 Pi message/event。
\end{itemize}
Agent loop 不应该出现 ``if provider === anthropic''。若新增 Provider 必须修改 loop，抽象边界就失败了。

\section{faux provider 为什么重要}
\pathref{packages/ai/src/providers/faux.ts} 不是玩具。它让测试精确配置模型响应、工具调用、token 切片、延迟、usage 和 prompt cache 模拟。真实 Provider 测试会受到密钥、账单、网络、模型升级和限流影响；核心 Agent 不变量必须优先由 faux provider 守护。

\section{三个 cache 不要混淆}
\begin{longtable}{p{0.25\textwidth}p{0.30\textwidth}p{0.35\textwidth}}
\toprule
概念 & 所在位置 & Pi 做什么\\
\midrule
模型 KV cache & GPU/推理服务器内部 attention tensor & Pi 不实现、不存储、不调度\\
Provider prompt cache & OpenAI/Anthropic 等服务端缓存稳定 prompt 前缀 & Pi 设置 session/cache 相关协议字段，统计 cache usage\\
本地 session/context & Pi JSONL、summary、session ID & Pi 管理恢复与稳定上下文，不是 tensor cache\\
\bottomrule
\end{longtable}

在 \pathref{packages/agent/src/agent.ts}，\texttt{sessionId} 被传给每轮流。不同 adapter 再将其映射为合适机制：OpenAI Responses 可使用 \texttt{prompt\_cache\_key}；Anthropic 可插入 \texttt{cache\_control} marker；部分后端使用 session affinity header。相关类型在 \pathref{packages/ai/src/types.ts}，OpenAI key 长度处理在 \pathref{api/openai-prompt-cache.ts}。

\begin{warningbox}
prompt cache 命中不等于 Pi 在本地保存了 KV tensor；它只表示后端可能复用稳定前缀。真正影响命中的工程行为包括保持 system prompt、工具定义和早期消息的稳定顺序，以及避免无意义地重写 session ID。
\end{warningbox}

\chapter{Extensions、skills 与 TUI}

\section{资源为何需要 trust}
项目目录可以携带 AGENTS.md、skills、prompt templates、themes、extension。它们可能改变 Agent 行为，extension 甚至能运行代码。因此 Pi 先做 project trust，再加载项目资源；不能先执行本地 extension 后再询问用户是否信任。

\section{extension 的边界}
extension 不应直接抓取 AgentSession 私有字段。Pi 通过 resource loader、extension runner、事件/hook/context 暴露能力。这样 lifecycle 可以记录和清理，错误可以诊断，session 切换时也能正确 teardown。

\section{TUI 是事件消费者}
\pathref{packages/tui} 是独立包，提供 terminal、editor、component、keybinding 和差量渲染。Coding Agent 的 \pathref{modes/interactive/interactive-mode.ts} 负责把 session event 映射为组件状态。

正确方向：
\begin{Verbatim}[fontsize=\small]
keyboard -> AgentSession.prompt / abort / steer
AgentSession event -> component state -> differential render
\end{Verbatim}
错误方向是让 UI 自行 append assistant 文本、执行工具或维护第二份 transcript；那会让 print、RPC、session restore 和 interactive 产生不同事实。

\chapter{如何系统学到可以维护 Pi}

\section{三条线并行}
\begin{itemize}
\item 60\%：当前代码纵向切片，理解今天怎样工作；
\item 25\%：mini Pi 复刻，证明能重新实现不变量；
\item 15\%：按需 Git 考古，解释设计为何变成今天这样。
\end{itemize}

\section{推荐顺序}
\begin{enumerate}
\item faux provider + 一轮文本流；
\item 单工具 loop；
\item Agent event/state 与 abort；
\item read tool 与输出截断；
\item AgentSession 和 JSONL session；
\item compaction；
\item bash 与安全；
\item Models/auth/一个真实 Provider；
\item extension/skill/trust；
\item pi-tui 与 interactive mode；
\item 每主题回溯四个关键 commit：引入、重构、bug fix、当前。
\end{enumerate}

\section{能力门槛}
\begin{longtable}{p{0.20\textwidth}p{0.70\textwidth}}
\toprule
门槛 & 证明\\
\midrule
最小闭环 & 能画 tool call 的 event/transcript 顺序并运行 faux test\\
安全修改 & 能为已有行为补测试或小修复，说明影响边界\\
产品重建 & 能从零实现 session、工具、context 与 print CLI\\
成熟设计解释 & 能用当前代码和四个 commit 解释 compaction/OAuth/TUI 等专题\\
贡献能力 & 能定位 issue、提出最小 patch、添加确定性 regression test\\
\bottomrule
\end{longtable}

\section{第一次读码作业}
\begin{enumerate}
\item 运行第 2 章的 faux-provider test；
\item 在 \pathref{packages/agent/test/e2e.test.ts} 找到该测试的两条 faux response；
\item 追 \texttt{Agent.prompt()}、\texttt{runLoop()}、\texttt{executeToolCalls()}；
\item 画出 user、assistant(tool call)、toolResult、assistant(final) 四条 transcript entry；
\item 写下：tool result 为什么要有 call ID？为什么 final assistant 不是第一次 assistant？
\end{enumerate}

\chapter{事件契约：把异步过程变成可观察事实}

\section{为什么事件模型比回调堆更重要}
一个 Agent 同时会发生很多事情：模型在输出文字；模型可能生成 tool call；工具正在运行；用户可能要求取消；UI 要及时显示；session 要可靠写盘。若每个模块随意直接调用另一个模块，就会产生隐形时间依赖：UI 先看到结果还是 session 先落盘？tool hook 在 assistant message 入历史前还是后执行？取消发生时哪个 promise 负责收尾？

Pi 的答案是将重要阶段建模为 \term{AgentEvent}。事件不是调试日志，而是跨层协议。它使每个消费者只订阅自己关心的事实：
\begin{itemize}
\item session listener 在 \texttt{message\_end} 写入最终消息；
\item TUI 在 \texttt{message\_update} 更新流式组件；
\item footer 在 usage 或 turn 结束时更新 token/cost；
\item extension hook 在规定时点检查工具调用；
\item 测试把事件序列作为行为规范。
\end{itemize}

\section{三个时间点}
理解 event 时，不要只问 ``发生了什么''，还要问 ``此刻 state 中有什么''。以 tool call 为例：
\begin{enumerate}
\item \texttt{message\_end(assistant)}：assistant 的完整 tool call 已进入 transcript；
\item \texttt{tool\_execution\_start}：参数已经解析并通过基本验证，\texttt{pendingToolCalls} 包含 call ID；
\item \texttt{tool\_execution\_end}：工具已成功或失败，pending 集合已移除该 ID；紧随其后才会有稳定的 tool result message。
\end{enumerate}

这也是为什么 \texttt{Agent.processEvents()} 先更新 state 再 await listener。若 listener 先执行，UI、session 和 extension 会观察到旧状态，出现难以复现的竞态。

\section{事件序列如何阅读}
把 \pathref{packages/agent/test/e2e.test.ts} 当作协议说明。它不只检查最终答案，也检查工具执行前后 \texttt{pendingToolCalls} 的集合。一个更完整的人工追踪表如下：

\begin{longtable}{p{0.25\textwidth}p{0.31\textwidth}p{0.31\textwidth}}
\toprule
事件 & transcript 状态 & 典型消费者\\
\midrule
\texttt{agent\_start} & 未必已有新增消息 & loading indicator、run timer\\
\texttt{message\_start(user)} & 用户消息正在进入 run & echo/输入确认\\
\texttt{message\_start(assistant)} & partial assistant 已建立 & assistant component\\
\texttt{message\_update} & partial 内容递增 & streaming markdown/thinking\\
\texttt{message\_end(assistant)} & final assistant 可持久化 & session、tool preflight\\
\texttt{tool\_execution\_start} & call 标记为 pending & tool card、取消状态\\
\texttt{tool\_execution\_update} & 中间进度可显示 & bash progress、长任务 UI\\
\texttt{message\_end(toolResult)} & 结果可成为下一轮 context & session、transcript view\\
\texttt{turn\_end} & 一个模型响应及其工具批次完成 & compaction/retry 决策\\
\texttt{agent\_end} & 本轮不再发事件 & final flush、idle UI\\
\bottomrule
\end{longtable}

\section{练习：从事件推状态}
假设你看到以下片段：
\begin{Verbatim}[fontsize=\small]
message_end(assistant with toolCall "edit")
tool_execution_start(edit, call-7)
tool_execution_end(edit, call-7, isError=true)
message_end(toolResult call-7)
turn_end
turn_start
\end{Verbatim}
请推断：第一轮是否已经结束？模型是否已经收到失败信息？为什么还有第二个 \texttt{turn\_start}？正确答案是：工具失败不是 run 失败；失败被标准化为 tool result，下一轮模型获得修正自己参数或策略的机会。

\chapter{错误、取消、重试：可靠 Agent 的负路径}

\section{成功路径只是最小部分}
演示 Agent 时常只展示 ``read 文件后回答''。真实使用中更常见的是网络断开、Provider 过载、无效工具参数、文件不存在、shell 退出非零、用户突然取消、context overflow。一个成熟 harness 的质量，更多由失败语义而不是 happy path 决定。

\section{四类失败}
\begin{longtable}{p{0.22\textwidth}p{0.31\textwidth}p{0.34\textwidth}}
\toprule
失败来源 & 例子 & 正确归宿\\
\midrule
输入/配置 & 没有模型、无凭据、坏 session ID & 启动诊断或明确 CLI error\\
Provider 流 & timeout、5xx、rate limit、stream error & final assistant error 或受控 auto-retry\\
工具执行 & 文件不存在、schema 不符、bash 失败 & \texttt{isError=true} tool result\\
用户控制 & abort、session 切换、关闭终端 & AbortSignal，最终状态为 aborted\\
\bottomrule
\end{longtable}

\section{abort 不是异常的别名}
\texttt{Agent.abort()} 触发当前 run 的 \texttt{AbortController}。该 signal 被交给 Provider stream 和 tool execute。abort 的语义是 ``用户不再需要这项工作''；它并不等同 ``系统出错''。因此 final assistant 的 \texttt{stopReason} 可以是 \texttt{aborted}，UI 可以显示已取消而不是红色错误，session 仍可以保留已经完成的历史。

对 bash 特别要警惕：取消上层 promise 不保证子进程已经停止。Coding Agent 为命令执行维护更具体的 process 管理逻辑；这是 product tool 层的职责，不应塞进通用 Agent loop。

\section{为什么不让工具返回一段 ``Error:'' 文本}
工具应在失败时抛出异常，由 Agent 统一转成 error tool result。好处是：
\begin{itemize}
\item 有单独的 \texttt{isError} 机器可读位，而不是依赖字符串前缀；
\item UI 可以以错误样式展示；
\item hook/metrics 能区分成功与失败；
\item 模型收到完整、稳定的失败消息而非半成功内容。
\end{itemize}

\section{retry 的边界}
自动重试只适用于暂时性 Provider 失败，例如 rate limit 或 5xx；它不应重试错误 tool 参数、权限拒绝或用户 abort。更重要的是重试不得重复执行有副作用的工具。Pi 将 Provider retry、tool execution、session compaction 分别处理，正是为了避免一个笼统 ``catch then retry'' 带来的重复写文件和重复命令。

\section{检查问题}
\begin{enumerate}
\item 如果文本已经 stream 到一半，用户按取消，session 中应保存什么？
\item 如果 read 工具失败，为什么下一轮仍值得调用模型？
\item 为什么 provider retry 不应该包住 ``模型调用 + 工具执行'' 的整个 loop？
\end{enumerate}

\chapter{系统提示词、上下文与 Agent 的工作记忆}

\section{prompt 不是一个字符串}
初学者常把系统提示词理解为一段固定文本。对 Coding Agent 而言，模型每轮请求的有效上下文至少包含：
\begin{itemize}
\item system prompt：角色、工具约束、行为规则；
\item 对话历史：用户、assistant、tool result；
\item 工具定义：名称、description、JSON schema；
\item 动态资源：项目 instruction、skills、prompt template；
\item Provider 选项：thinking、cache、session affinity、transport。
\end{itemize}
它们共同决定模型能看到什么、能调用什么、后续请求能否复用缓存。

\section{稳定前缀为何重要}
很多 Provider 的 prompt cache 依赖稳定前缀。若每一轮随机重排 system prompt、工具定义或早期 context，后端很难复用已处理部分；成本、延迟和 cache read 都会变差。Pi 的 session ID、cache retention 和 context 构造因此不是微小性能优化，而是 Agent 长对话体验的一部分。

\section{transformContext 的正确位置}
\texttt{transformContext} 在 AgentMessage 层工作，发生在转换为 LLM Message 之前。它适合：
\begin{enumerate}
\item 裁剪过长工具输出；
\item 注入显式 summary；
\item 筛掉仅供 UI 使用的 entry；
\item 加入外部但可信的项目上下文。
\end{enumerate}
它不应该偷偷删掉用户历史而不留下可审计痕迹。正确做法是 compaction：将旧片段转为 summary entry，并令用户、session 和模型都能解释这个转换。

\section{上下文窗口是产品约束}
模型的 context window 并不是越大越不用管理。长上下文仍有成本、延迟、注意力稀释和 Provider 限制；工具输出会快速放大 transcript。Pi 的 compaction、truncate、token estimate、cache statistics 都是在处理这个产品约束。学习时把 ``上下文预算'' 当作和内存预算同样真实的资源。

\section{练习：画一份请求快照}
任选一条测试，在模型调用点写下：
\begin{Verbatim}[fontsize=\small]
systemPrompt:
messages:
tools:
model:
sessionId:
thinking level:
cache options:
\end{Verbatim}
下一轮在 append tool result 后再次写一份。比较两份快照，便能看清 Agent loop 为什么不是普通聊天。

\chapter{测试策略：用确定性模拟代替昂贵的猜测}

\section{测试是最短设计文档}
Pi 的测试价值不在覆盖率数字，而在把异步行为固定为可读的契约。比起从 production code 猜 ``并行工具完成后结果如何排序''，先看 \pathref{packages/agent/test/agent-loop.test.ts} 的断言会更快得到答案。

\section{四层测试}
\begin{enumerate}
\item \term{协议单测}：schema、消息转换、cache option、路径/截断 helper。
\item \term{loop 单测}：给定 scripted assistant message，断言 event 与 transcript。
\item \term{session/product 测试}：JSONL、恢复、fork、compaction、settings、tools。
\item \term{受控交互测试}：TUI、CLI/RPC mode；少量真实 Provider smoke test。
\end{enumerate}
越靠内层，越应确定、快速、无需凭据。越靠外层，越少而精。

\section{faux provider 的能力}
\pathref{providers/faux.ts} 可设置一串 response，或提供根据 context 返回 message 的 factory。它可以模拟：
\begin{itemize}
\item text/thinking/tool-call content blocks；
\item 每秒 token 数和 token chunk 大小；
\item usage 与 cache read/write；
\item 第二轮是否正确看到第一轮 tool result；
\item abort 时是否产生正确 final message。
\end{itemize}

这让你能写一种真正有价值的测试：``给定 assistant 先要求 calculate，再给最终文本，Agent 应生成哪些状态与事件？'' 而不是 ``请真实模型算 2+2，希望它合作''。

\section{测试一个工具闭环}
理想的断言至少包括：
\begin{Verbatim}[fontsize=\small]
assistant message has toolCall(call-1)
tool_execution_start(call-1)
toolResult.toolCallId === call-1
toolResult.isError === false
next assistant message sees result
pendingToolCalls is empty after completion
\end{Verbatim}
若只断言最终自然语言答案，即使 tool 根本没执行、模型碰巧猜对了，也会测试通过。

\section{读测试的顺序}
\begin{enumerate}
\item 先读 test name：作者想保证什么？
\item 再读 setup：faux model 和 tools 被如何配置？
\item 列出断言的事件、state 和消息；
\item 最后跳到实现中找使断言成立的最短函数路径。
\end{enumerate}
这种 ``test -> implementation'' 路径比先读几千行 runtime 更适合初学复杂 Agent。

\chapter{用 Git 历史理解设计，而不是淹没在 Git 历史里}

\section{何时应该查历史}
当前 Pi 有数千个 commit。线性阅读会混入版本 bump、模型目录生成、格式化和 CI 维护，学习效率很低。应只在以下问题出现时查 history：
\begin{itemize}
\item 这个看似复杂的分层为什么存在？
\item 某个兼容字段或额外 event 在防哪类 bug？
\item 为什么这里选择并发/串行、summary/删除、session ID/随机 ID？
\item 哪个测试是回归测试，它最初修了什么？
\end{itemize}

\section{每个主题四个 commit}
对一个文件或概念，最多选择：
\begin{enumerate}
\item 初次引入：需求最原始的形状；
\item 关键重构：职责为何重新划分；
\item 高价值 bug fix：不变量曾如何被打破；
\item 当前版本：最终形态。
\end{enumerate}

例如 tool execution 可以围绕 \pathref{agent-loop.ts} 查并行批次、完成顺序、truncated tool call 拒绝、per-tool sequential override。每次都比较父提交和目标提交，先写 ``变化前什么会失败''，再读 diff。

\section{安全的 Git 命令}
\begin{Verbatim}[fontsize=\small]
git log --follow -- packages/agent/src/agent-loop.ts
git log --follow -- packages/coding-agent/src/core/agent-session.ts
git show --stat <commit>
git diff <commit>^ <commit> -- <path>
git show <commit>:<path>
\end{Verbatim}
不要为了阅读历史频繁把主工作目录切到旧 commit。使用 \texttt{git show}、临时 worktree 或只读 detached checkout；实验代码放独立目录。

\section{从 history 得到什么}
历史不是为了记 SHA，而是为了捕获设计动机。例如 ``state before listener'' 这种看似微小的顺序，往往来自实际 race bug；它比注释更能说明为什么不能随意重排。把这种结论写入你的 invariants 笔记，才是考古的产物。

\chapter{从 Pi 扩展到更完整的 AI Agent 版图}

\section{Pi 已覆盖什么}
Pi 是学习单 Agent engineering 的高密度案例：
\begin{itemize}
\item 流式 LLM 协议统一；
\item structured tool calling；
\item context/token/cache usage；
\item session、compaction、retry；
\item 多 Provider、认证和 OAuth；
\item coding tool 安全、containerization；
\item event-driven terminal UX。
\end{itemize}

\section{Pi 不主打什么}
以下不是 Pi 的主要实现重心：
\begin{itemize}
\item 多 Agent 协作、planner/worker/supervisor 图；
\item RAG、向量数据库、retrieval ranking；
\item 模型训练、推理引擎、GPU 显存/KV cache 管理；
\item 企业级长期记忆、评估平台和 observability backend。
\end{itemize}
这不是缺陷，而是边界清晰。理解 Pi 后，补 LangGraph 学图式编排，补 RAG 学检索，补 vLLM/TGI 学推理服务，补 eval/trace 系统学可靠性，路径会更自然。

\section{一个概念映射}
\begin{longtable}{p{0.29\textwidth}p{0.28\textwidth}p{0.29\textwidth}}
\toprule
问题 & Pi 的答案 & 可后续学习的相邻领域\\
\midrule
单 Agent 如何反复调用工具？ & Agent loop + tool result & ReAct、function calling\\
多个步骤如何显式分叉？ & 有限的 queue/session 分支 & LangGraph/workflow engines\\
如何让模型知道私有资料？ & project context/skills，不是检索系统 & RAG/vector search\\
如何降低重复长 prompt 成本？ & Provider prompt cache/session ID & inference cache/KV cache\\
如何隔离危险命令？ & 外部 sandbox/containerization & policy sandbox/VM orchestration\\
\bottomrule
\end{longtable}

\chapter{八周学习路线与学习日志模板}

\section{第 1--2 周：核心闭环}
\begin{itemize}
\item 跑 help、model list、faux provider e2e；
\item 读 \texttt{Agent.prompt}、\texttt{runLoop}、\texttt{streamAssistantResponse}；
\item 画 user -> assistant(tool) -> toolResult -> assistant 的四消息图；
\item 写 mini loop，只支持一个 echo/read tool。
\end{itemize}

\section{第 3--4 周：产品状态}
\begin{itemize}
\item 读 AgentSession、session manager、read tool；
\item 实现 JSONL append/reload 和路径受限 read；
\item 学 abort、error tool result、parallel/sequential；
\item 为 mini loop 增加 faux tests。
\end{itemize}

\section{第 5--6 周：长对话与 Provider}
\begin{itemize}
\item 读 compaction、token estimate、cache usage；
\item 选一个 OpenAI 或 Anthropic adapter，不同时读多个；
\item 区分 prompt cache 和 KV cache；
\item 实现显式 summary，不静默删历史。
\end{itemize}

\section{第 7--8 周：扩展与交互}
\begin{itemize}
\item 读 resource loader、trust、extension runner；
\item 独立读 pi-tui 的 component/render contract；
\item 最后进入 interactive mode；
\item 选择一个主题回溯四个关键 commit，写设计说明。
\end{itemize}

\section{每次学习日志}
\begin{Verbatim}[fontsize=\small]
场景：
运行命令与可观察输出：
测试名称：
入口 -> 核心 -> 输出：
状态由谁拥有：
消息/事件不变量：
失败、取消、重试路径：
一个相关 commit 与它解决的问题：
我的最小实验：
仍不理解的问题：
\end{Verbatim}
当这些日志积累到十几篇，你将拥有比一份泛泛教程更有价值的个人 Pi 设计文档。

\chapter{源码地图：从问题定位到正确文件}

\section{不要用 ``全仓搜索'' 代替架构}
Pi 的当前实现中，同一个词可能出现在 Provider、Agent、Coding Agent、TUI 和测试多个层。搜索 ``session'' 会产生大量结果，却不告诉你哪一层拥有 session；搜索 ``tool'' 更会混合 schema、执行、渲染、extension 和文档。先按问题选择包，再用符号和测试缩小范围。

\section{问题到包的第一映射}
\begin{longtable}{p{0.30\textwidth}p{0.29\textwidth}p{0.29\textwidth}}
\toprule
你正在问什么？ & 第一站 & 第二站\\
\midrule
命令行 flag 为何如此处理？ & \texttt{coding-agent/src/cli/args.ts} & \texttt{main.ts}\\
为什么进入 print/RPC/interactive？ & \texttt{main.ts} & \texttt{modes/}\\
模型为何收到某些消息？ & \texttt{agent-loop.ts} & \texttt{types.ts}/context transform\\
工具为何执行/被拒绝/并发？ & \texttt{agent-loop.ts} & \texttt{coding-agent/core/tools/}\\
状态何时写入 session？ & \texttt{agent-session.ts} & \texttt{session-manager.ts}\\
某 Provider 请求长什么样？ & \texttt{ai/src/api/<adapter>.ts} & \texttt{models.ts}/auth\\
为什么 cache usage 变化？ & \texttt{ai/src/types.ts} & Provider adapter/faux test\\
某个终端组件为何刷新？ & \texttt{tui/src/} & \texttt{interactive-mode.ts}\\
extension 能做什么？ & \texttt{core/extensions/} & resource loader/trust manager\\
\bottomrule
\end{longtable}

\section{核心文件的阅读顺序}
\begin{enumerate}
\item \pathref{packages/agent/src/types.ts}：先定义词汇。不要猜 \texttt{AgentTool}、\texttt{AgentEvent}、\texttt{AgentContext} 的字段。
\item \pathref{packages/agent/src/agent.ts}：看 state、public API、listener 屏障、AbortController、queues。
\item \pathref{packages/agent/src/agent-loop.ts}：看每一轮的实际控制流。
\item \pathref{packages/ai/src/types.ts} 与 \pathref{models.ts}：看模型流、Provider、auth 的契约。
\item \pathref{packages/coding-agent/src/core/agent-session.ts}：看产品策略如何附着在 core 上。
\item \pathref{packages/coding-agent/src/core/tools/}：看真实 side effect 如何被工具化。
\end{enumerate}

每看一个实现文件，都优先找同名测试。例如 \texttt{agent-loop.ts} 对应 \texttt{agent-loop.test.ts}；不要先试图理解 coding agent 的巨型 interactive mode。

\section{一个实际定位例子}
问题：``为什么两个 tool 同时执行时，界面先显示第二个完成，但历史仍是第一个在前？''
\begin{enumerate}
\item 这是 Agent runtime 时间顺序问题，不是 TUI 问题，先到 \texttt{packages/agent}；
\item 搜索 \texttt{tool\_execution\_end}，找到 \texttt{agent-loop.test.ts} 的 completion/source order 测试；
\item 读测试的 setup 和断言，确认两种顺序的预期；
\item 再读 \texttt{executeToolCallsParallel()}，看 ``completion event eager'' 与 ``persisted result source order'' 的实现位置；
\item 最后才看 TUI 如何展示 end event。
\end{enumerate}
这个过程说明：好的源码导航从\term{不变量}开始，不从文件名或 UI 现象开始。

\section{命令工具箱}
\begin{Verbatim}[fontsize=\small]
# 当前实现中某符号的定义/使用
rg "executeToolCallsParallel" packages/agent

# 某文件的演进史
git log --follow -- packages/agent/src/agent-loop.ts

# 某个旧版本文件，不切换工作树
git show <commit>:packages/agent/src/agent-loop.ts

# 当前文件和一个 commit 前后的差异
git diff <commit>^ <commit> -- packages/agent/src/agent-loop.ts
\end{Verbatim}
搜索结果只是假设；测试和调用链才是证据。

\chapter{Mini Pi 实验室：用复刻证明你真的理解}

\section{实验室的角色}
mini Pi 不是竞争项目，也不应复制 upstream 的实现。它是\term{理解校验器}：每次你读懂 Pi 一项机制，就在几十到几百行教学代码中重建它的关键不变量。若自己写不出，或写出的行为无法由测试固定，说明需要回到 Pi 的测试和类型。

\section{建议的最小数据模型}
\begin{Verbatim}[fontsize=\small]
type UserMessage = { role: "user"; content: TextPart[]; timestamp: number };
type AssistantMessage = {
  role: "assistant";
  content: Array<TextPart | ToolCallPart>;
  stopReason: "stop" | "toolUse" | "error" | "aborted";
};
type ToolResultMessage = {
  role: "toolResult";
  toolCallId: string;
  toolName: string;
  content: TextPart[];
  isError: boolean;
};
\end{Verbatim}
先只支持 text 和 tool call。不要一开始支持图片、thinking、多个 Provider 或 extension。每增加一种 content block，都要重新问：如何流式更新？如何持久化？如何转换为 Provider 负载？如何在 UI 中显示？

\section{阶段一：scripted model}
定义一个不联网的 model：
\begin{Verbatim}[fontsize=\small]
interface Model {
  stream(
    messages: AgentMessage[],
    tools: AgentTool[],
    signal?: AbortSignal
  ): AsyncIterable<AssistantDelta>;
}
\end{Verbatim}
第一条 scripted response 返回 ``hello''；第二个版本先返回 \texttt{read} tool call，检查输入里是否已有 tool result，再返回总结。这样你可以证明 loop 确实发起了第二次模型调用，而不是应用代码偷偷拼接结果。

\section{阶段二：工具和错误}
给 tool 设置统一 execute 契约：
\begin{Verbatim}[fontsize=\small]
type AgentTool = {
  name: string;
  execute(args: unknown, signal?: AbortSignal):
    Promise<{ content: TextPart[] }>;
};
\end{Verbatim}
未知工具、坏参数、读不到文件都要产生 \texttt{isError=true} result。不要在 tool 内把异常吞成空字符串。为每个错误写一个测试：模型能否看到 error result？Agent 是否仍能继续第二轮？

\section{阶段三：并发与稳定历史}
实现两个 read tool call 并行。故意让第一个慢、第二个快。你的实现应允许第二个先发 ``完成'' UI event，但 transcript 中 tool result 仍按 assistant 源顺序排列。若只用 \texttt{Promise.all} 然后立即按 resolve 顺序 append，你就复现不了 Pi 的不变量。

\section{阶段四：session 与 compaction}
以 JSONL append message；重启后读回并重建 context。然后加入一个非常简单的 ``超过 N 条消息就产生 summary'' 规则。你不需要先让模型生成精美摘要；关键是 summary 必须显式进入历史，且测试能说明哪几条原消息被替代。

\section{mini Pi 的验收表}
\begin{longtable}{p{0.35\textwidth}p{0.54\textwidth}}
\toprule
功能 & 最小验收\\
\midrule
text stream & partial 不写入最终 transcript\\
tool call & tool result ID 与 call ID 一致\\
tool error & 错误不使整个 Agent 无提示崩溃\\
parallel tools & completion 与 persistence 顺序分别可断言\\
abort & final stop reason 是 aborted，后续工具不再启动\\
session & reload 后下一次模型看到相同上下文\\
summary & 压缩是可见 entry，不是悄悄删除\\
\bottomrule
\end{longtable}

\chapter{Provider adapter 解剖：把不同模型协议收敛为一个流}

\section{为什么 adapter 是必要复杂度}
OpenAI、Anthropic、Google、Bedrock 等服务的请求字段、tool 格式、thinking、SSE event、cache control 和 usage 名称各不相同。如果 Agent loop 直接面向它们，循环会被大量 provider 条件污染，新增模型也会改变核心行为。Pi 将这种复杂度局限在 \pathref{packages/ai/src/api/}。

\section{adapter 的输入与输出}
一个 adapter 的输入包含 Pi 的 \texttt{Model}、\texttt{Context}、标准 stream options；输出是 \texttt{AssistantMessageEventStream}。它要完成四类转换：
\begin{enumerate}
\item Pi system/messages/tools $\rightarrow$ Provider request；
\item Provider stream events $\rightarrow$ Pi text/thinking/tool-call delta；
\item Provider final usage/stop reason $\rightarrow$ Pi final assistant message；
\item Pi session/cache/auth options $\rightarrow$ Provider-specific headers/body fields。
\end{enumerate}

\section{先读 faux，再读一个真实 adapter}
\pathref{providers/faux.ts} 比真实 adapter 更适合先读，因为它完整展示 Pi 侧事件、usage 和 context 的目标形状，没有 SDK 噪声。理解 faux 后，只选择一个真实 adapter：
\begin{itemize}
\item 选择 \pathref{api/openai-responses.ts}，学习 Responses API、prompt cache key；
\item 或选择 \pathref{api/anthropic-messages.ts}，学习 content blocks、thinking、cache control。
\end{itemize}
不要同时读两者。先写一张 ``相同 Pi 消息如何映射'' 的对照表，再扩展到另一个。

\section{认证不应泄漏到 Agent loop}
\texttt{Models} 负责把 credential store、环境变量、OAuth token 或 Provider-specific auth 解析为请求选项。这样 Agent 只关心 ``可否 stream''，而不是 ``token 是否过期''。这一边界也使 status UI 能准确显示 ``未配置''、``需要重新登录'' 与 ``Provider 网络失败'' 的区别。

\section{Adapter 阅读问题}
\begin{enumerate}
\item 此 Provider 如何编码 tool definition 和 tool result？
\item 哪些 Provider event 被合并为同一种 Pi delta？
\item provider stop reason 如何变成 Pi 的 stop reason？
\item usage 的 cache read/write 如何归一？
\item session ID 在请求中是 header、body key，还是不使用？
\item 出错时 stream 是否仍能给 Agent 一个 final result？
\end{enumerate}
若这些问题能回答，你读的不是 SDK 调用，而是在理解 Pi 的跨 Provider 契约。

\chapter{终端 UI 深读前需要懂的渲染机制}

\section{为什么终端 UI 仍是一个独立系统}
LLM 每秒可生成多次 delta。若每个 delta 都清屏并重画整个聊天记录，终端会闪烁、慢、破坏 scrollback；复杂 terminal 还会出现光标错乱。因此 \pathref{packages/tui} 不是几段 ANSI escape 的集合，而是组件树、终端抽象、输入 buffer、键绑定与差量渲染系统。

\section{差量渲染的抽象}
\begin{Verbatim}[fontsize=\small]
Agent event
  -> component state changes
  -> component tree renders terminal lines
  -> compare with previous frame
  -> emit minimal ANSI updates
\end{Verbatim}
这和浏览器 DOM diff 类似，但约束不同：终端宽度会变化、Unicode 宽度不恒定、scrollback 不能乱改、输入编辑器需要光标位置、某些 terminal 不支持所有 capability。

\section{Pi 的两层 UI}
\begin{itemize}
\item \pathref{packages/tui/src/}：通用组件，如 text、markdown、editor、select list、terminal；
\item \pathref{packages/coding-agent/src/modes/interactive/}：Pi 产品组件，如 assistant message、tool execution、footer、model/session selector。
\end{itemize}
先学通用 TUI 的 render contract，再读 interactive mode 把 \texttt{AgentSessionEvent} 转成什么组件。倒过来读会淹没在 import 和产品细节中。

\section{输入的正确方向}
输入并不是直接修改 transcript：
\begin{Verbatim}[fontsize=\small]
keyboard -> editor -> command/input handler
         -> AgentSession.prompt / steer / abort
         -> AgentSession event -> rendered component
\end{Verbatim}
这保证同一 AgentSession 可被 print/RPC 使用，且 session restore 后 UI 只是重放事实，不是重新推断聊天内容。

\section{第一次 TUI 作业}
不要启动完整 interactive mode 后盲目阅读。先找一个简单 \texttt{Text} 或 \texttt{Markdown} component，理解它如何输出 lines；再找 \texttt{TUI} 如何启动/触发 render；最后追一个 \texttt{message\_update} 在 interactive mode 中如何更新 assistant component。完成这一小条链之后，再研究 editor、overlay、keybinding。

\chapter{术语表、检查清单与下一次阅读的入口}

\section{核心术语表}
\begin{longtable}{p{0.25\textwidth}p{0.64\textwidth}}
\toprule
术语 & 在 Pi 中的准确含义\\
\midrule
Agent harness & 为模型调用、工具、状态、事件和上下文提供运行时边界的基础设施；不是某个特定 prompt。\\
turn & 一次模型响应及紧随其后的工具批次。一个用户任务可包含多个 turn。\\
transcript & Agent/session 记录的消息历史。它比单纯 LLM context 更宽，可含 custom/summary entry。\\
context & 发送给模型的 system prompt、转换后的 messages 和 tools。它从 transcript 构建，不必与 transcript 一致。\\
partial message & 还在 stream 中的 assistant 状态，用于 UI；不能当作最终持久化事实。\\
final message & stream 完成后经 \texttt{result()} 得到的完整 assistant 消息，含最终 usage、stop reason、工具参数。\\
tool call & assistant content block 中的结构化调用意图，包含 name、arguments 和唯一 call ID。\\
tool result & 工具执行后的标准消息，必须关联 call ID，并显式标记是否 error。\\
preflight & 工具真正执行前的查找、参数校验、policy/hook 检查阶段。\\
execution mode & 一个 tool batch 的串行或并行策略；副作用工具可要求整个 batch 顺序执行。\\
steering & 正在运行时输入的优先新指令；当前工具批次完成后进入下一 turn。\\
follow-up & Agent 原本结束后才处理的排队工作。\\
AbortSignal & 取消 stream/tool 的协作信号；不是普通 error。\\
faux provider & Pi 的确定性假模型，用于不花钱、不联网的 Agent 测试。\\
Provider & 拥有模型列表、认证语义和实际流实现的供应商适配单元。\\
Models & Provider 集合与认证/委托门面；Agent 通过它间接请求模型。\\
session ID & 本地 session、Provider cache key 或 session affinity 的关联标识，具体映射由 adapter 决定。\\
prompt cache & 服务端复用稳定 prompt 前缀的机制；与本地 JSONL 和 GPU KV cache 不同。\\
KV cache & 推理服务内部保存 attention key/value tensor 的机制；Pi 不直接管理。\\
compaction & 将旧上下文压缩为显式 summary，以满足 context budget 的过程。\\
\bottomrule
\end{longtable}

\section{第一次完整调用链的空白模板}
每次阅读一个功能，都复制下面模板，不要只写 ``看了哪些文件''：
\begin{Verbatim}[fontsize=\small]
用户场景：
触发入口（CLI / RPC / UI）：
模式选择：
AgentSession 如何创建：
Agent API：
传入模型的 Context：
Provider/stream 事件：
Agent event：
工具调用和结果：
session 写入点：
UI/print 输出点：
失败与取消路径：
相关测试：
一个关键历史 commit：
\end{Verbatim}

\section{阅读 Agent loop 前的自测}
若以下任一问题回答不清，不要急着深入 \texttt{agent-loop.ts} 的 helper：
\begin{enumerate}
\item user message 何时进入 \texttt{agent.state.messages}？
\item \texttt{message\_update} 中消息为什么不能直接写 session？
\item assistant 返回 tool call 后，为什么还要再调用模型？
\item tool result 的 call ID 为什么不能省略？
\item tool throw 和 Provider stream error 分别变成什么？
\item abort、steering、follow-up 各自在哪个时间点生效？
\item 并发 tool 的完成顺序为何不等于 transcript 顺序？
\end{enumerate}

\section{阅读 Provider 前的自测}
\begin{enumerate}
\item 为什么 Agent loop 不应该 import OpenAI 或 Anthropic SDK？
\item \texttt{Provider} 和 \texttt{Models} 的职责如何区分？
\item 为什么 faux provider 对核心测试比真实 API 更可靠？
\item prompt cache、session ID、session JSONL 和 GPU KV cache 的边界是什么？
\item 新增 Provider 时，哪个层应该改变，哪个层不应该改变？
\end{enumerate}

\section{阅读 Coding Agent 前的自测}
\begin{enumerate}
\item 为什么 \texttt{AgentSession} 不能简单等同 \texttt{Agent}？
\item 为什么 cwd 改变要重建 services？
\item read/write/bash 分别有哪些不属于 Agent core 的产品风险？
\item 为什么 project resource 加载要经过 trust？
\item 为什么 interactive UI 不应修改第二份对话历史？
\end{enumerate}

\section{建议的学习记录评分法}
每完成一个纵向切片，自评 0--2 分：
\begin{longtable}{p{0.48\textwidth}p{0.15\textwidth}p{0.25\textwidth}}
\toprule
能力 & 分数 & 说明\\
\midrule
可以复现一个测试/场景 & 0--2 & 0=未运行；1=运行但不懂；2=能解释 setup/断言\\
可以画调用链 & 0--2 & 0=只有文件名；1=大致路径；2=有状态和输出边界\\
可以说明不变量 & 0--2 & 0=只复述代码；1=知道结果；2=知道为何不能破坏\\
可以预测失败路径 & 0--2 & 0=不清楚；1=知道一种 error；2=区分 error/abort/retry\\
可以写最小实验 & 0--2 & 0=无；1=改了但无测试；2=有确定性验证\\
可以解释历史动机 & 0--2 & 0=未查；1=看过 diff；2=能说前后不变量变化\\
\bottomrule
\end{longtable}
总分低于 8 分时，不要进入下一大主题；回到测试或 mini Pi。总分 10 分以上，才值得加入一个新层，例如 session、cache 或 TUI。

\section{下一次打开源码时，从这里开始}
第一条完整学习任务只需这些文件：
\begin{Verbatim}[fontsize=\small]
packages/agent/test/e2e.test.ts
packages/agent/src/agent.ts
packages/agent/src/agent-loop.ts
packages/agent/src/types.ts
packages/ai/src/providers/faux.ts
\end{Verbatim}
先运行本书第 2 章命令，再从测试中找到 ``executes tools and tracks pending tool calls''。不要额外打开 interactive mode、真实 Provider 或 compaction。等你能独立画出它的四消息闭环，再移动到 \pathref{packages/coding-agent/src/core/agent-session.ts}。

\chapter{第一天的 30 分钟任务}

\section{目标}
今天不要求理解 Pi，只要建立第一个不依赖 API key 的完整观察。结束时你应该能明确说出：faux model 在哪里设定了两轮响应，工具何时执行，结果如何返回模型。

\section{时间盒}
\begin{enumerate}
\item \textbf{第 0--5 分钟：}打开 \pathref{packages/agent/test/e2e.test.ts}，只找到测试名 \texttt{executes tools and tracks pending tool calls}。先不读实现。
\item \textbf{第 5--10 分钟：}运行本书第 2 章的 Vitest 命令。确认只有该测试执行；记录 ``1 passed''。
\item \textbf{第 10--15 分钟：}读该测试中 \texttt{faux.setResponses([...])} 的两条 response。第一条包含什么 content block？第二条为什么是最终文本？
\item \textbf{第 15--20 分钟：}读 \pathref{packages/agent/test/utils/calculate.ts}，确认 tool 接受什么参数、如何生成结果。
\item \textbf{第 20--25 分钟：}在 \pathref{packages/agent/src/agent.ts} 找到 \texttt{prompt()}，在 \pathref{agent-loop.ts} 找到 \texttt{executeToolCalls()}。只看函数入口与调用关系。
\item \textbf{第 25--30 分钟：}用下面模板写五行笔记。
\end{enumerate}

\begin{Verbatim}[fontsize=\small]
用户消息：
第一次 assistant message：
工具调用（名称、call ID、参数）：
tool result：
第二次 assistant message：
\end{Verbatim}

\section{今天不要做什么}
不要配置真实 Provider；不要打开 \texttt{interactive-mode.ts}；不要尝试理解 compaction；不要改生产代码。第一次成功的标准不是修改了什么，而是你能在纸上画出消息闭环，并知道哪条测试能证明它。

\backmatter
\chapter*{下一版会补什么}
\addcontentsline{toc}{chapter}{下一版会补什么}
v1.0 会在本书结构上增加：Agent loop 的逐段源码阅读、AgentSession services/runtime、JSONL/fork/compaction 数据模型、read/edit/write/bash 的实现细节、一个真实 Provider adapter、prompt cache 与 session affinity、extension lifecycle、TUI 差量渲染，以及每个主题四个关键 commit 的演进说明。核心版已经建立这些深读的地图；下一版不会推翻本书，而是沿每章继续向下钻。

\end{document}
